\section{Foundations of Causal Discovery}
\label{sec:foundations}

Causal discovery infers causal relationships from observational data. Constraint-based
methods, such as PC and FCI, use conditional-independence tests. PC returns a
Markov equivalence class of DAGs, while FCI also allows latent confounding and
returns a partial ancestral graph \cite{spirtes2000causation}.

Score-based methods search for a DAG that fits the data well. GES performs a greedy
search over equivalence classes. NOTEARS reformulates DAG learning as continuous
optimization by expressing acyclicity as a smooth constraint
\cite{zheng2018dags}. DYNOTEARS extends this formulation to contemporaneous and
lagged relationships in time series \cite{pamfil2020dynotears}.

Functional causal models make stronger assumptions about the functional form and
noise distribution. LiNGAM can identify the causal order when the model is linear
and the noise is non-Gaussian \cite{shimizu2006lingam}. VAR-LiNGAM extends this
idea to vector autoregressive models with instantaneous and lagged effects
\cite{hyvarinen2010estimation}.

These foundational methods generally assume one causal model for the full dataset.
Real data may combine several latent processes, populations, time intervals, or
biological states. A single model fitted to such data can produce an averaged graph
that represents none of the underlying mechanisms. This motivates methods for
heterogeneous structure learning.

\subsection{Causal Discovery from Nonstationary and Heterogeneous Data}
\label{ssec:nonstationary}

Several methods use changes across environments to improve causal identification.
Invariant Causal Prediction identifies parents whose conditional relationship with a
target remains stable across known environments \cite{peters2016causal}. Its main
limitation is the need for environment labels.

CD-NOD uses distribution shifts across domains or time to recover causal structure
\cite{huang2020causal}. It needs an auxiliary domain or time index, but not hard
environment labels. CD-NOTS adapts this idea to time series and improves the
treatment of lagged variables \cite{sadeghi2025causal}. The Sparse Mechanism Shift
hypothesis gives theoretical support to this approach when only a few mechanisms
change \cite{perry2022causal}.

DyCAST models continuous changes in the causal DAG with Neural ODEs
\cite{cheng2025dycast}. IRM provides a related learning-theoretic approach that
encourages stable predictors across environments \cite{arjovsky2019irm}. Both IRM
and ICP require known environment partitions.

Together, these studies show that heterogeneity can help causal learning. However,
existing methods generally need environment labels, an auxiliary nonstationarity
index, or temporal ordering. Recovering regimes without all three is the central gap
addressed by this thesis.

\subsection{Mixture Models and Causal Clustering}
\label{ssec:clustering}

Causal clustering jointly estimates cluster assignments and causal structures. This
is preferable to a sequential approach because incorrect assignments can corrupt
graph learning, while an incorrect graph can corrupt clustering.

The Specific and Shared Causal Model (SSCM) represents data as a mixture of causal
mechanisms and separates shared from regime-specific relations
\cite{huang2019specific}. It requires the number of regimes in advance.

CCSL is more flexible. It uses a causal Chinese Restaurant Process to group
observations by causal-structure similarity rather than feature distance
\cite{chen2023ccsl}. Variational inference estimates the parameters and assignments
together, and the number of groups need not be supplied. CCSL provides
identifiability under a linear non-Gaussian model and has been tested on fMRI and
Sachs protein-signalling data.

\subsubsection*{i. Temporal Extensions of Causal Clustering}

MCD learns mixtures of structural causal models from time series using variational
inference \cite{varambally2024mcd}. Its autoregressive structure makes it unsuitable
for i.i.d. cross-sectional data. CASTOR provides a temporal method for jointly
learning sequential regimes and their DAGs \cite{rahmani2025castor}, but its use of
temporal contiguity does not transfer directly to unordered observations.

\subsubsection*{ii. i.i.d. Extensions and Recent Advances}

HCL jointly learns latent clusters and causal structures from mixed observational
data without environment labels or a pre-specified regime count
\cite{li2025hcl}. Its preprint reports weaker assumptions than CCSL, but its
identifiability conditions and large-graph performance require further study.
MCVCI and MCVCC show that mechanism-aware clustering is also useful in simple
bivariate mixtures \cite{liu2025hybrid}.

\subsubsection*{iii. Classical Clustering Methods}

Methods such as k-means, GMM, DBSCAN, and OPTICS group observations by similarity
in feature space. They can fail when regimes have similar distributions but
different causal mechanisms. CCSL results show the value of including causal
structure in the clustering objective \cite{chen2023ccsl}.

\subsection{Joint Causal Inference from Multiple Contexts}

JCI represents contexts as additional nodes in a causal graph and learns causal
structure across observational and interventional settings \cite{mooij2020jci}.
It is general, but the contexts must be known in advance. Representation-learning
methods also study hidden causal variables, but they usually focus on learning
latent variables from images or videos rather than grouping tabular observations by
causal regime.
